\documentclass{umk-matematika}

\begin{document}

\maketitlepage
    {Практическая работа №26}
    {Сфера и шар}
    {}

\section{Фундамент (1--4 балла)}

\textbf{Задача 1 (1 балл).} Радиус шара $5$ см. Найдите площадь его поверхности.

\textbf{Задача 2 (2 балла).} Радиус шара $3$ см. Найдите его объём.

\textbf{Задача 3 (3 балла).} Площадь поверхности шара $64\pi$ см². Найдите радиус.

\textbf{Задача 4 (4 балла).} Объём шара $288\pi$ см³. Найдите радиус.

\section{Стены (5--7 баллов)}

\textbf{Задача 5 (5 баллов).} Сечение шара плоскостью, удалённой от центра на $8$ см, имеет площадь $36\pi$ см². Найдите радиус шара.

\textbf{Задача 6 (6 баллов).} Радиус шара $13$ см. Сечение находится на расстоянии $5$ см от центра. Найдите площадь сечения.

\textbf{Задача 7 (7 баллов).} Два шара имеют радиусы $3$ см и $6$ см. Во сколько раз объём большего шара больше объёма меньшего?

\section{Кровля (8--10 баллов)}

\textbf{Задача 8 (8 баллов).} Резервуар имеет форму шара радиусом $5$ м. Сколько краски потребуется для покраски его поверхности, если расход краски $0{,}2$ кг на 1 м²?

\textbf{Задача 9 (9 баллов).} Купол здания имеет форму полусферы радиусом $10$ м. Сколько краски потребуется для покраски внешней поверхности купола, если расход краски $0{,}2$ кг на 1 м²?

\textbf{Задача 10 (10 баллов).} Шар радиусом $15$ см переплавлен в цилиндр, радиус основания которого равен радиусу шара. Найдите высоту цилиндра.

\newpage
\section*{Ответы}

\begin{enumerate}
  \item $100\pi$ см²
  \item $36\pi$ см³
  \item $4$ см
  \item $6$ см
  \item $10$ см
  \item $144\pi$ см²
  \item в $8$ раз
  \item $20\pi$ кг $\approx 62{,}8$ кг
  \item $40\pi$ кг $\approx 125{,}6$ кг
  \item $20$ см
\end{enumerate}

\end{document}